\chapter{BiCG and BiCGSTAB}\label{v0ch:14}
\begin{chaptersummary}
Because GMRES basis storage and orthogonalization grow rapidly on large nonsymmetric systems, short-recurrence alternatives are often attractive. This chapter develops BiCG through two-sided Krylov spaces and BiCGSTAB through local residual smoothing. Short recurrences sharply cut memory usage while introducing nonmonotone convergence and potential breakdown, making solver selection a balance among memory, matrix structure, and robustness.
\end{chaptersummary}
\learninggoals{Interpret shadow residuals and biorthogonality; execute BiCGSTAB; distinguish breakdown from convergence.}
\paragraph{Core Question}Can nonsymmetric systems retain short recurrences, and at what cost?

\section{Two-Sided Projection}
BiCG constructs $\K_k(A,\bm r_0)$ alongside $\K_k(A^T,\widetilde{\bm r}_0)$, enforcing biorthogonal residuals and biconjugate search directions with a shadow residual satisfying $\widetilde{\bm r}_0^T\bm r_0\ne0$. Because this requires products with both $A$ and $A^T$ and key inner products can vanish while the true residual remains large, BiCGSTAB instead multiplies the BiCG residual polynomial by local linear smoothing factors, eliminating explicit $A^T$ applications and damping erratic oscillations, though residual decrease remains nonmonotone.
\section{An Executable Recurrence}
Starting from a fixed shadow vector $\widetilde{\bm r}$, initialize $\bm p=\bm v=0$ and
$\rho_{\rm old}=\alpha=\omega=1$. Each iteration evaluates
\[
 \rho=\widetilde{\bm r}^T\bm r,\quad
 \beta=(\rho/\rho_{\rm old})(\alpha/\omega),\quad
 \bm p=\bm r+\beta(\bm p-\omega\bm v),\quad
 \bm v=A\bm p,\quad \alpha=\rho/(\widetilde{\bm r}^T\bm v),
\]
\[
 \bm s=\bm r-\alpha\bm v,\qquad
 \bm t=A\bm s,\qquad
 \omega=\frac{\bm t^T\bm s}{\bm t^T\bm t},\qquad
 \bm x\leftarrow\bm x+\alpha\bm p+\omega\bm s,\quad
 \bm r\leftarrow\bm s-\omega\bm t.
\]
If $\bm s$ already satisfies the stopping tolerance, updating by $\alpha\bm p$ and exiting avoids a $0/0$ division; otherwise, monitoring $\rho$, $\widetilde{\bm r}^T\bm v$, $\bm t^T\bm t$, and $\omega$ detects breakdown. Under right preconditioning, both search directions are transformed by $M^{-1}$ and the transformed vectors enter the solution update directly.
\conceptbridge{What does the short recurrence save, and what does it give up?}{GMRES retains its entire Krylov basis to minimize the residual globally over the current subspace. BiCGSTAB keeps only a constant number of vectors via coupled polynomial recurrences, trading away stepwise residual optimality and introducing sensitivity to near-zero inner products.}{Full-basis minimization $\longrightarrow$ short-recurrence selection $\longrightarrow$ storage savings and breakdown risks.}
\section{A Minimal Experiment}
On a one-dimensional upwind advection--diffusion tridiagonal matrix with exact solution $\bm x_*=\bm1$ and $\bm b=A\bm x_*$, comparing GMRES and BiCGSTAB by true-residual history, operator evaluations, and wall time—while varying diffusion strength and right-hand sides—reveals how nonnormality and initial residuals shape convergence.
\costnote{Each step performs two operator and preconditioner applications with $O(N)$ storage, so total runtime reflects both per-step work and convergence rate. When breakdown or stagnation occurs, inspecting the return status and true residual guides preconditioner tuning, shadow-vector restarts, or switching to GMRES.}
\coderef{sparse_linear_algebra/solvers.py}
\aiexercise{For $A=\left[\begin{smallmatrix}3&1\\0&2\end{smallmatrix}\right]$ and $\bm b=(1,1)^T$, calculate the first step and identify the local residual minimized by $\omega$.}{Construct a near-breakdown example. Use scale-aware denominator checks and distinguish early convergence from genuine failure.}
\studyroute{Full two-sided BiCG derivation and stabilized-polynomial variants are omitted. See \cite{saad,templates}.}{BiCG; biorthogonality; shadow residual; BiCGSTAB; breakdown}
